%! suppress = EscapeUnderscore
\section{Basics}

\bd[Machine Learning (Arthur Samuel)]
\textbf{Machine learning} is the field of study that gives computers the ability to learn without being explicitly
programmed.
\ed

\bd[Machine Learning (Tom Mitchell)]
A computer program is said to learn from experience $E$ with respect to a task $T$ and a performance measure $P$, if its
performance on $T$, as measured by $P$, improves with experience $E$.
\ed

Machine learning is the study of computer algorithms that can improve automatically through experience and by the use
of data. It is seen as a part of artificial intelligence. Machine learning algorithms build a model based on sample
data, known as training data, in order to make predictions or decisions without being explicitly programmed to do so.
Machine learning algorithms are used in a wide variety of applications, such as in medicine, email filtering,
speech recognition, and computer vision, where it is difficult or unfeasible to develop conventional algorithms to
perform the needed tasks.

\subsection{Use Cases}

Machine learning has found increasing usage in both enterprise and consumer applications. As its adoption in the
industry quickly grows, machine learning has proven to be a powerful tool for a wide range of problems. Despite an
incredible amount of excitement and hype generated by people both inside and outside the field, machine learning is
not a magic tool that can solve all problems. Even for problems that machine learning can solve, machine learning
solutions might not be the optimal solutions. Before starting a machine learning project, you might want to ask
whether machine learning is necessary or cost-effective. \v

The essence of machine learning is that a pattern exists, and it can not be pined down mathematically, however we have
data on it, and we can treat it in a probabilistic way. Thus, machine learning is great for:
\bit
\item \textbf{Systems that have the capacity to learn}: For a machine learning system to learn, there must be something
for it to learn from. In most cases, machine learning systems learn from data.
\item \textbf{Data is available, or it's possible to collect data}: Because machine learning learns from data, there
must be data for it to learn from.
\item \textbf{There are patterns to learn, and they are complex}: Complex problems for which using a traditional
approach yields no good solution, the best machine learning techniques can perhaps find a solution.
\item \textbf{Existing solutions require a long lists of rules}: One machine learning algorithm can often simplify code
and perform better than the traditional approach.
\item \textbf{It's a predictive problem}: Machine learning models make predictions, so they can only solve problems that
require predictive answers.
\item \textbf{Unseen data shares patterns with the training data}: The patterns your model learns from existing data are
only useful if unseen data also share these patterns. In technical terms, it means your unseen data and training data
should come from similar distributions.
\item \textbf{It's repetitive}: When a task is repetitive, each pattern is repeated multiple times, which makes it
easier for machines to learn it. Despite exciting progress in few-shot learning research, most machine learning
algorithms still require many examples to learn a pattern.
\item \textbf{The cost of wrong predictions is cheap}: Unless your machine learning model's performance is 100\% all the
time, which is highly unlikely for any meaningful tasks, your model is going to make mistakes. Machine learning is
especially suitable when the cost of a wrong prediction is low.
\item \textbf{It's at scale}: Machine learning solutions often require nontrivial up-front investment on data, compute,
infrastructure, and talent, so it'd make sense if we can use these solutions a lot. ``At scale'' means different things
for different tasks, but, in general, it means making a lot of predictions. A problem might appear to be a singular
prediction, but it's actually a series of predictions.
\item \textbf{The patterns are constantly changing}: If your problem involves one or more constantly changing patterns,
hardcoded solutions such as handwritten rules can become outdated quickly. Figuring how your problem has changed so that
you can update your handwritten rules accordingly can be too expensive or impossible. Because machine learning models
learn from data,you can update your machine learning model with new data without having to figure out how the data has
changed.
\eit

Some usual examples of machine learning algorithms are: analyzing images to automatically classify them, detecting
tumors in brain scans, automatically classifying news articles, automatically flagging offensive comments on
discussion forums, creating a chat-bot or a personal assistant, forecasting company revenue, making an application
reacting to voice commands, detecting credit card fraud, segmenting clients based on their purchases, representing a
complex high-dimensional dataset in a clear and insightful way, recommending a product that a client may be
interested in, building an intelligent bot for a game, translating from one language to another, personal assistants,
and many more \ldots \v

According to Algorithmia's 2020 state of enterprise machine learning survey, machine learning applications in 
enterprises are diverse, serving both internal use cases (reducing costs, generating customer insights and intelligence, 
internal processing automation) and external use cases (improving customer experience, retaining customers, interacting 
with customers) as shown in the following figure.

\fig{ml_use_cases}{0.5}

\subsection{Requirements}

We can't say that we've successfully built a machine learning system without knowing what requirements the system has
to satisfy. The specified requirements for an machine learning system vary from use case to use case. However, most
systems should have the following four characteristics:
\bit
\item \textbf{Reliability}: The system should continue to perform the correct function at the desired level of
performance even in the face of adversity (hardware or software faults, and even human error).
\item \textbf{Scalability}: Whichever way your system grows (complexity, traffic volume, machine learning models),
there should be reasonable ways of dealing with that growth, both in resource scaling, but also in artifact management.
\item \textbf{Maintainability}: It's important to structure your workloads and set up your infrastructure in such a
way that different contributors can work using tools that they are comfortable with, instead of one group of
contributors forcing their tools onto other groups.
\item \textbf{Adaptability}: Because machine learning systems are part code, part data, and data can change quickly,
machine learning systems need to be able to evolve quickly. To adapt to shifting data distributions and business
requirements, the system should have some capacity for both discovering aspects for performance improvement and
allowing updates without service interruption.
\eit

\subsection{Academia VS Industry}

As machine learning usage in the industry is still fairly new, most people with machine learning expertise have gained
it through academia: taking courses, doing research, reading academic papers. However, machine learning in production
(industry) is very different from machine learning in research (academia). Here follows a list with some important
differences.
\bit
\item \textbf{Different Stakeholders \& Requirements}: People involved in a research project often align on one single
objective. The most common objective is model performance—develop a model that achieves the state-of-the-art results on
benchmark datasets. To edge out a small improvement in performance, researchers often resort to techniques that make
models too complex to be useful. On the other hand, there are many stakeholders involved in bringing a machine learning
system into production. Each stakeholder has their own requirements. Having different, often conflicting, requirements
can make it difficult to design, develop, and select a machine learning model that satisfies all the requirements.
\item \textbf{Computational Priorities}: When designing a machine learning system, people who haven't deployed a machine
learning system often make the mistake of focusing too much on the model development part and not enough on the model
deployment and maintenance part. During the model development process, you might train many different models, and each
model does multiple passes over the training data. Each trained model then generates predictions on the validation data
once to report the scores. The validation data is usually much smaller than the training data. During model development,
training is the bottleneck. Once the model has been deployed, however, its job is to generate predictions, so inference
is the bottleneck. Research usually prioritizes fast training, whereas production usually prioritizes fast inference.
One corollary of this is that research prioritizes high throughput whereas production prioritizes low latency.
\item \textbf{Data}: During the research phase, the datasets you work with are often clean and well-formatted, freeing
you to focus on developing models. They are static by nature so that the community can use them to benchmark new
architectures and techniques. This means that many people might have used and discussed the same datasets, and quirks of
the dataset are known. You might even find open source scripts to process and feed the data directly into your models.
In production, data, if available, is a lot more messy. It's noisy, possibly unstructured, constantly shifting. It's
likely biased, and you likely don't know how it's biased. Labels, if there are any, might be sparse, imbalanced, or
incorrect. Changing project or business requirements might require updating some or all of your existing labels. If
you work with users' data, you'll also have to worry about privacy and regulatory concerns. In research, you mostly work
with historical data, e.g., data that already exists and is stored somewhere. In production, most likely you'll also
have to work with data that is being constantly generated by users, systems, and third-party data.
\item \textbf{Fairness}: During the research phase, a model is not yet used on people, so it's easy for researchers to
put off fairness as an afterthought. When it gets to production, it's too late. If you optimize your models for better
accuracy or lower latency, you can show that your models beat state of the art. But there's no equivalent state of the
art for fairness metrics. Machine learning algorithms don't predict the future, but encode the past, thus perpetuating
the biases in the data and more. When machine learning algorithms are deployed at scale, they can discriminate against
people at scale. If a human operator might only make sweeping judgments about a few individuals at a time, a machine
learning algorithm can make sweeping judgments about millions in split seconds. This can especially hurt members of
minority groups because misclassification on them could only have a minor effect on models' overall performance metrics.
\item \textbf{Interpretability}: Since most machine learning research is still evaluated on a single objective, model
performance, researchers aren't incentivized to work on model interpretability. However, interpretability isn't just
optional for most machine learning use cases in the industry, but a requirement. First, interpretability is important
for users, both business leaders and end users, to understand why a decision is made so that they can trust a model and
detect potential biases mentioned previously. Second, it's important for developers to be able to debug and improve a
model.
\eit

Let's summarize these bullet points in the following figure.

\fig{ml_industry_vs_academia}{0.65}

\subsubsection{Objectives}

Another big difference between academia and industry is the matter of objectives. In any machine learning project
there are two sets of objectives. Machine learning and business objectives. Starting with the objectives of the
proposed machine learning projects, when working on a machine learning project, data scientists tend to care about
the machine learning objectives: the metrics they can measure about the performance of their machine learning systems.
On the other hand, most companies don't care about machine learning metrics but about business metrics. Since the
sole purpose of businesses is to maximize profits for shareholders, the ultimate goal of any project within a
business is, therefore, to increase profits. For a machine learning project to succeed within a business organization,
it's crucial to tie the performance of a machine learning system to the overall business performance. \v

In reality, the effect of a machine learning project on business objectives can be hard to reason about. To gain a
definite answer on the question of how machine learning metrics influence business metrics, experiments are often needed.
Many companies do that with experiments like A/B testing and choose the model that leads to better business metrics,
regardless of whether this model has better machine learning metrics. \v

Yet, even rigorous experiments might not be sufficient to understand the relationship between a machine learning
model's outputs and business metrics. When evaluating machine learning solutions through the business lens, it's
important to be realistic about the expected returns. Due to all the hype surrounding machine learning, generated
both by the media and by practitioners with a vested interest in machine learning adoption, some companies might have
the notion that machine learning can magically transform their businesses overnight. There are many companies that
have seen payoffs from machine learning, but this gain hardly happened overnight. In general returns on investment in
machine learning depend a lot on the maturity stage of adoption. The longer you've adopted machine learning, the more
efficient your pipeline will run, the faster your development cycle will be, the less engineering time you'll need,
and the lower your cloud bills will be, which all lead to higher returns. \v 

In general, some good questions to ask while framing the problem and defining the objectives in business terms are:
\bit
\item How will your solution be used?
\item What are the current solutions/workarounds (if any)?
\item How should you frame this problem?
\item How should performance be measured?
\item Is the performance measure aligned with the business objective?
\item What would be the minimum performance needed to reach the business objective?
\item What are comparable problems? Can you reuse experience or tools?
\item Is human expertise available? How would you solve the problem manually?
\eit

Last but not least, make sure to list the assumptions you (or others) have made so far and verify them if possible.

\subsubsection{Mind VS Data}

Progress in the last decade shows that the success of an machine learning system depends largely on the data it was
trained on. Instead of focusing on improving machine learning algorithms, most companies focus on managing and
improving their data. Despite the success of models using massive amounts of data, many are skeptical of the emphasis
on data as the way forward. This is the main argument of many public debates on the power of mind versus data. \v

Mind might be disguised as inductive biases or intelligent architectural designs. Data might be grouped together with
computation since more data tends to require more computation. In theory, you can both pursue architectural designs
and leverage large data and computation, but spending time on one often takes time away from another. The debate isn't
about whether finite data is necessary, but whether it's sufficient. The term finite here is important, because if
we had infinite data, it might be possible for us to look up the answer. Having a lot of data is different from having
infinite data. \v

Regardless of which camp will prove to be right eventually, no one can deny that data is essential, for now. Both the
research and industry trends in the recent decades show the success of machine learning relies more and more on the
quality and quantity of data. Models are getting bigger and using more data.

\subsection{Machine Learning VS Software Engineering}

Since machine learning is part of software engineering, and software has been successfully used in production for more
than half a century, some might wonder why we don't just take tried-and-true best practices in software engineering and
apply them to machine learning. That's an excellent idea. In fact, machine learning production would be a much better
place if machine learning experts were better software engineers. Many traditional software engineering tools can be
used to develop and deploy machine learning applications.

\fig{ml_vs_swe}{0.8}

\bit
\item \textbf{Code \& Data}: In software engineering, there's an underlying assumption that code and data are
separated. In fact, in software engineering, we want to keep things as modular and separate as possible. On the
contrary, machine learning systems are part code, part data, and part artifacts created from the two. The trend in
the last decade shows that applications developed with the most/best data win. Instead of focusing on improving
machine learning algorithms, most companies will focus on improving their data. Because data can change quickly,
machine learning applications need to be adaptive to the changing environment, which might require faster development
and deployment cycles.
\item \textbf{Testing \& Versioning}: In traditional software engineering, you only need to focus on testing and
versioning your code. With machine learning, we have to test and version our data too, and that's the hard part. How
to version large datasets? How to know if a data sample is good or bad for your system? Not all data samples are
equal—some are more valuable to your model than others. Indiscriminately accepting all available data might hurt your
model's performance and even make it susceptible to data poisoning attacks.
\item \textbf{Size}: The size of machine learning models is another challenge. As of 2022, it's common for machine
learning models to have hundreds of millions, if not billions, of parameters, which requires gigabytes of RAM to load
them into memory. There is the question of how to get these models to run fast enough to be useful.
\item \textbf{Monitoring \& Debugging}: Monitoring and debugging these models in production is also nontrivial. As
machine learning models get more complex, coupled with the lack of visibility into their work, it's hard to figure out
what went wrong or be alerted quickly enough when things go wrong.
\eit

\subsection{Types}

There are so many different types of machine learning systems that it is useful to classify them in broad categories,
based on the following criteria:
\ben
\item Whether or not they are trained with human supervision. Based on this category we have the following
subcategories:
\bit
\item \textbf{Supervised Learning}: In supervised learning the training set you feed to the algorithm includes the
desired solutions, called labels. Some of the most important supervised learning algorithms are: Linear Regression,
Logistic Regression, Support Vector Machines (SVMs), k-Nearest Neighbors, Decision Trees, Random Forests, Neural
Networks.
\item \textbf{Unsupervised Learning}: In unsupervised learning the training data is unlabeled so the system tries to
learn without a teacher. Some of the most important unsupervised learning algorithms are:
\bit
\item Clustering algorithms like K-Means, DBSCAN and Hierarchical Cluster Analysis (HCA).
\item Anomaly detection algorithms such as One-Class SVM and Isolation Forest.
\item Dimensionality reduction algorithms such as Principal Component Analysis (PCA), Kernel PCA, Locally Linear
Embedding (LLE), and t-Distributed Stochastic Neighbor Embedding (t-SNE).
\eit
\item \textbf{Semisupervised Learning}: In semisupervised learning one has plenty of unlabeled instances, and few
labeled instances. Most semisupervised learning algorithms are combinations of unsupervised and supervised algorithms.
\item \textbf{Reinforcement Learning}: In reinforcement learning the learning system, called an agent, can observe
the environment, select and perform actions, and get rewards in return. It must then learn by itself what is the
best strategy, called a policy, to get the most reward over time. A policy defines what action the agent should
choose when it is in a given situation.
\eit
\item Whether or not they can learn incrementally on the fly. Based on this category we have the following
subcategories:
\bit
\item \textbf{Offline or Batch Learning}: In offline, or batch, leaning the system is incapable of learning
incrementally and it must be trained using all the available data. This will generally take a lot of time and
computing resources, so it is typically done offline. First the system is trained, and then it is launched into
production and runs without learning anymore, it just applies what it has learned.
\item \textbf{Online Learning}: In online learning one trains the system incrementally by feeding it data instances
sequentially, either individually or in small groups called ``mini-batches''. Each learning step is fast and cheap, so
the system can learn about new data on the fly, as it arrives.
\eit
\item Whether they work by simply comparing new data points to known data
points, or instead by detecting patterns in the training data and building a predictive model. Based on this
category we have the following subcategories:
\bit
\item \textbf{Instance-Based Learning}: In instance-based learning the system learns the examples by heart, then
generalizes to new cases by using a similarity measure to compare them to the learned examples (or a subset of them).
\item \textbf{Model-Based Learning}: In model-based learning one in order to generalize from a set of examples they
build a model of these examples and then use that model to make predictions.
\eit
\een

These criteria are not exclusive. You can combine them in any way you like. This fact makes machine learning is a
very broad topic with many different branches and applications. In these notes we will cover the vast majority of them.

\subsection{Development}

Developing a machine learning system is an iterative and, in most cases, never-ending process. Once a system is put
into production, it'll need to be continually monitored and updated. Here is a simplified representation of what the
iterative process for developing machine learning systems in production.

\bit
\item \textbf{Step 1. Project Scoping}: A project starts with scoping the project, laying out goals, objectives, and
constraints. Stakeholders should be identified and involved.
\item \textbf{Step 2. Data Engineering}: A vast majority of machine learning models today learn from data, so developing
machine learning models starts with engineering data.
\item \textbf{Step 3. Model Development}: With the initial set of training data, we'll need to extract features and
develop initial models leveraging these features. This is the stage that requires the most machine learning knowledge.
\item \textbf{Step 4. Model Deployment}: After a model is developed, it needs to be made accessible to users. Developing
a machine learning system is like writing—you will never reach the point when your system is done. But you do reach the
point when you have to put your system out there.
\item \textbf{Step 5. Monitoring \& Improvement}: Once in production, models need to be monitored for performance decay
and maintained to be adaptive to changing environments and requirements.
\item \textbf{Step 6. Business Analysis}: Model performance needs to be evaluated against business goals and analyzed to
generate business insights. These insights can then be used to eliminate unproductive projects or scope out new projects.
This step is closely related to the first step.
\eit

\fig{steps}{0.42}

\subsection{Conventions}

Before we end this section we have to mention that there are some common conventions in the machine learning 
community around the notation used to describe various notions. We will of course follow the same conventions. In 
order to briefly formalize the essence of machine learning we will introduce some of the very basic notation that we 
will be using throughout the notes now, although we will introduce more notation in the later chapters. Here are some
very basic concepts with their usual notation:
\bit
\item Input: $x \in X$.
\item Output: $y \in Y$.
\item Data: $\{ x_{i}, y_{i} \}, \:\:\: i=1,2,3,\ldots, m$.
\item Target Function: $f: X \to Y$.
\item Hypothesis Function: $h: X \to Y$ with $ h \approx f$.
\item Hypothesis Set: $H = \{h\}$.
\eit

Having defined these mathematical notation we can now say that, informally, the goal of machine learning is, based on
the data $\{ x_{i}, y_{i} \}$, to discover a hypothesis function $h$, out of a set of possible hypothesis functions $H$,
that behaves in a similar way with the target function $f$ which is, and always will be, unknown to us.

\fig{img/mlmodel}{0.39}

The question is how can we learn an unknown function $f$ just based on the data we already have, when the unknown
function $f$ in general can take any value outside the known data. The short answer is that we can not however,
without proving it, the following relation holds:
\bse
P \Big[ | E_{\text{in}} (h) - E_{\text{out}} (h) | > \epsilon \Big] \leq 2 \cdot M \cdot e^{2\epsilon^2m}
\ese

where $ E_{\text{in}} (h)$ is the error that we get for $h$ in the known data, $E_{\text{out}} (h)$ is the error that
we will get when we use $h$ for new data, $M$ is the number of possible hypothesis function $h$ (i.e.\ the cardinality
of the hypothesis set $H = \{h\}$, $\epsilon$ is the tolerance that we have for errors, and $m$ is the number of data
points. This equation tells us that no matter what, learning is possible only in a probabilist sense. We will always
have an error, since the whole process carries a stochastic nature. \v

So we can informally summarize what we are trying to do with machine learning as:
\bit
\item From aforementioned relation: $E_{\text{in}} \approx E_{\text{out}}$.
\item From learning algorithm: $E_{\text{in}} \approx 0$.
\item From the combination of these 2: $E_{\text{out}} \approx 0$.
\eit

By having $E_{\text{out}} \approx 0$, that means that our hypothesis function $h$ generalizes well for out of sample
data, so we can use it for predictions. That in a nutshell is how machine learning works.

\section{Data}

As we have described in the previous section, machine learning is tightly coupled with data. A machine learning
system can work with data from many different sources. They have different characteristics, can be used for different
purposes, and require different processing methods. Understanding the sources your data comes from can help you use
your data more efficiently.

\subsection{Data Sources}

\bd[User Input Data]
\textbf{User input data} are data explicitly input by users.
\ed

User input data can be text, images, videos, uploaded files, etc, and can be easily malformatted. User input data
requires more heavy-duty checking and processing. On top of that, users also have little patience. In most cases,
when we input data, we expect to get results back immediately. Therefore, user input data tends to require fast
processing.

\bd[System Generated Data]
\textbf{System generated data} are data generated by different components of your systems.
\ed

System generated data include various types of logs (which record the state and significant events of the system),
and system outputs such as model predictions. Because these data are system generated, they are much less likely to
be malformatted the way user input data is. Overall, they don't need to be processed as soon as they arrive, the way
you would want to process user input data.

\bd[Internal Databases]
\textbf{Internal databases} are generated by various services and enterprise applications in a company.
\ed

Internal databases manage assets such as inventory, customer relationship, users, and more. This kind of data can be
used by machine learning models directly or by various components of an machine learning system. The data inside
internal databases are usually called ``first party data''.

\bd[First Party Data]
\textbf{First party data} is the data that your company already collects about your users or customers.
\ed

Similarly there are ``second party data'' and ``third party data''.

\bd[Second Party Data]
\textbf{Second party data} is the data collected by another company on their own customers that they make available
to you, though you'll probably have to pay for it.
\ed

\bd[Third Party Data]
\textbf{Third party data} is the data that companies collect on the public who aren't their direct customers.
\ed

\subsection{Data Formats}

Once you have data, you might want to store it (or ``persist'' it, in technical terms). Since your data comes from
multiple sources with different access patterns, storing your data isn't always straightforward and, for some cases,
can be costly. It's important to think about how the data will be used in the future so that the format you use
will make sense. \v

The process of converting a data structure or object state into a format that can be stored or transmitted and
reconstructed later is called ``serialization''.

\bd[Serialization]
\textbf{Serialization} is the process of translating a data structure or object state into a format that can be stored
(for example, in a file or memory data buffer) or transmitted (for example, over a computer network) and reconstructed
later (possibly in a different computer environment).
\ed

There are many data serialization formats. When considering a format to work with, you might want to consider
different characteristics such as human readability, access patterns, and whether it's based on text or binary, which
influences the size of its files. For more on serialization check the ``Serialization'' section on ``Technologies''
chapter in ``Computer Science'' part.

\subsubsection{Structured \& Unstructured Data}

\bd[Structured Data]
\textbf{Structured data} is data that has a standardized format for efficient access by software and humans.
\ed

\bd[Data Warehouse]
A repository for storing structured data is called a \textbf{data warehouse}.
\ed

Data warehouses are used to store data that has been processed into formats ready to be used. \v

Structured data follows a predefined data model, also known as a ``database schema''.

\bd[Data Schema]
The \textbf{data schema} is the structure of a database described in a formal language supported by the database
management system. The term ``schema'' refers to the organization of data as a blueprint of how the database is
constructed.
\ed

\be
For example, the database schema might specify that each data item consists of two values: the first value, \code{name},
is a string of at most 50 characters, and the second value, \code{age}, is an 8-bit integer in the range between 0 and
200. The predefined structure makes your data easier to analyze. If you want to know the average age of people in the
database, all you have to do is to extract all the age values and average them out.
\ee

The disadvantage of structured data is that you have to commit your data to a predefined schema. If your schema
changes, you'll have to retrospectively update all your data, often causing mysterious bugs in the process. Because
business requirements change over time, committing to a predefined data schema can become too restricting. Or you
might have data from multiple data sources that are beyond your control, and it's impossible to make them follow the
same schema. This is where ``unstructured data'' becomes appealing.

\bd[Unstructured Data]
\textbf{Unstructured data} is data that either does not have a predefined data model or is not organized in a predefined
database schema.
\ed

\bd[Data Lake]
A repository for storing unstructured data is called a \textbf{data lake}.
\ed

Data lakes are usually used to store raw data before processing. \v

Unstructured data doesn't adhere to a predefined data schema. It's usually text but can also be numbers, dates,
images, audio, etc. Even though unstructured data doesn't adhere to a schema, it might still contain intrinsic
patterns that help you extract structures.

The following table shows a summary of the key differences between structured and unstructured data.

\fig{str_unst}{0.7}

\subsubsection{Extract, Transform, Load (ETL)}

In the early days of the relational data model, data was mostly structured. When data is extracted from different
sources, it's first transformed into the desired format before being loaded into the target destination such as a
database or a data warehouse.

\bd[Extract]
\textbf{Extract} is extracting the data you want from all your data sources.
\ed

Some of them will be corrupted or malformatted. In the extracting phase, you need to validate your data and reject
the data that doesn't meet your requirements. For rejected data, you might have to notify the sources. Since this is
the first step of the process, doing it correctly can save you a lot of time downstream.

\bd[Transform]
\textbf{Transform} is the part of the process, where most of the data processing is done.
\ed

During transform, you might want to join data from multiple sources and clean it. You might want to standardize the
value ranges. You can apply operations such as transposing, de-duplicating, sorting, aggregating, deriving new
features, more data validating, etc.

\bd[Load]
\textbf{Load} is deciding how and how often to load your transformed data into the target destination, which can be a
file, a database, or a data warehouse.
\ed

This process is called ``ETL'', which stands for ``extract, transform, load''.

\bd[Extract, Transform, Load (ETL)]
\textbf{Extract, transform, load} (\textbf{ETL}), is a three-phase process where data is extracted, transformed and
loaded into an output data container.
\ed

Even before machine learning, ETL was all the rage in the data world, and it's still relevant today for machine
learning applications. ETL refers to the general purpose processing and aggregating of data into the shape and the
format that you want. The idea of ETL sounds simple but powerful, and it's the underlying structure of the data layer
at many organizations. An overview of the ETL process is shown in the following figure.

\fig{etl}{0.35}

\subsection{Data Models}

\bd[Data Model]
A \textbf{data model} is an abstract model that organizes elements of data and standardizes how they relate to one
another and to the properties of real-world entities.
\ed

Data models describe how data is represented. How you choose to represent data not only affects the way your systems 
are built, but also the problems your systems can solve. In this section, we'll study two types of models that seem 
opposite to each other but are actually converging: relational models and NoSQL models. We'll go over examples to 
show the types of problems each model is suited for.

\subsubsection{Relational Model}

Relational models are among the most persistent ideas in computer science. Invented by Edgar F. Codd in 1970, the
relational model is still going strong today, even getting more popular.

\bd[Relational Model]
The \textbf{relational model} is an approach to managing data using a structure and language consistent with first-order
predicate logic where all data is represented in terms of tuples, grouped into relations.
\ed

The idea is simple but powerful. In this model, data is organized into relations; each relation is a set of tuples. A
table is an accepted visual representation of a relation, and each row of a table makes up a tuple. Relations are
unordered. You can shuffle the order of the rows or the order of the columns in a relation and it's still the same
relation. Data following the relational model is usually stored in file formats like \code{CSV}. \v

Databases built around the relational data model are relational databases.

\bd[Relational Database]
A database organized in terms of the relational model is called a \textbf{relational database}
\ed

Once you've put data in your databases, you'll want a way to retrieve it. The language that you can use to specify
the data that you want from a database is called a query language.

\bd[Query Language]
A \textbf{query language} is a computer language used to make queries in databases and information systems.
\ed

The most popular query language for relational databases today is SQL.

\bd[Structured Query Language (SQL)]
\textbf{Structured query language} (\textbf{SQL}) is a domain-specific query language used in programming and designed
for managing data held in a relational databases.
\ed

Even though inspired by the relational model, the data model behind SQL has deviated from the original relational
model. The most important thing to note about SQL is that it's a declarative language, as opposed to Python, which is
an imperative language.\footnote{In the imperative paradigm, you specify the steps needed for an action and the
computer executes these steps to return the outputs. In the declarative paradigm, you specify the outputs you want,
and the computer figures out the steps needed to get you the queried outputs.} With SQL, you specify the pattern of
data you want, the tables you want the data from, the conditions the results must meet, the basic data
transformations such, and many more, but not how to retrieve the data. It is up to the database system to decide how
to break the query into different parts, what methods to use to execute each part of the query, and the order in
which different parts of the query should be executed. In practice, it's not always easy to write a query to solve a
specific task, and it's not always feasible or tractable to execute a query.

\subsubsection{NoSQL}

The relational data model has been able to generalize to a lot of use cases, from ecommerce to finance to social
networks. However, for certain use cases, this model can be restrictive. For example, it demands that your data
follows a strict schema, and schema management is painful. It can also be difficult to write and execute SQL queries
for specialized applications. The latest movement against the relational data model is `NoSQL`.

\bd[NoSQL]
A \textbf{NoSQL} (originally referring to \textbf{non-SQL} or \textbf{non-relational}) database provides a mechanism for
storage and retrieval of data that is modeled in means other than the tabular relations used in relational databases.
\ed

Two major types of NoSQL models are the document model and the graph model. The document model targets use
cases where data comes in self-contained documents and relationships between one document and another are rare. The
graph model goes in the opposite direction, targeting use cases where relationships between data items are common and
important. Since NoSQL is quite uncommon in the field of machine learning, we will not go into any more detail.

\subsection{Get, Explore, Prepare...}

Switching from theory to practice, before we finish this introductory chapter, we will provide some useful and practical
tips about what to do with data in your every day life. We will split them into 3 categories: ``Get'', ``Explore'' and
``Prepare'' the data.

\subsubsection*{Get The Data}
\bit
\item List the data you need and how much you need.
\item Find and document where you can get that data.
\item Check how much space it will take.
\item Check legal obligations, and get access authorization if necessary.
\item Create a workspace with enough storage space.
\item Get the data.
\item Make sure you data are not of insufficient quantity, poor quality, or non representative of the population.
Especially validation and test set must be as representative as possible.
\item Convert the data to a format you can easily manipulate, without changing the data itself.
\item Ensure sensitive information is deleted or protected.
\item Make a copy of the original dataset and work with this one.
\item Automate as much as possible so you can easily get fresh data.
\eit

\subsubsection*{Explore The Data}
\bit
\item Create a Jupyter notebook to keep a record of your data exploration.
\item Study each attribute and its characteristics: name, type, percentage of missing values, noisiness and type of
noise, outliers, usefulness for the task, etc.
\item Study basic statistics of each attribute: mean, standard deviation, distribution, etc.
\item For supervised learning tasks, identify the targets.
\item Sample a test set, put it aside, and never look at it.
\item Visualize the data.
\item Study the correlations between attributes.
\item Study how you would solve the problem manually.
\item Identify extra data that would be useful.
\item Document what you have learned.
\item Try to get insights from a field expert for these steps.
\eit

\subsubsection*{Prepare The Data}
Write functions for all data transformations you apply, for the following reasons:
\bit
\item So you can easily prepare the data the next time you get a fresh dataset.
\item So you can apply these transformations in future projects.
\item To clean and prepare the test set.
\item To clean and prepare new data instances once your solution is live.
\item To make it easy to treat your preparation choices as hyperparameters.
\eit

\section{Models}

\subsection{Development}

The process of developing a machine learning model typically starts with a business objective\footnote{It's important
to recognize that ML projects don't happen in a vacuum. They are generally part of a larger project that in turn
impacts technologies, processes, and people. That means part of establishing objectives also includes change
management, which may even provide some guidance for how the ML model should be built.}. Business objectives
naturally come with performance targets, technical infrastructure requirements, and cost constraints; all of these
factors can be captured as key performance indicators (KPIs), which will ultimately enable the business performance
of models in production to be monitored. \v

With clear business objectives defined, it is time to bring together subject matter experts and data scientists to
begin the journey of developing the ML model. This starts with the search for suitable  data, which we covered in the
previous section. \v

After data collection the next step is training. Usually before developing any kind of model, machine learning
practitioners spend some time in the so called ``before machine learning'' phase. In other words, they try to make a
prediction with non-machine learning solutions with the simplest heuristics. Then they move to the second phase of
``developing the simplest machine learning models''. \v

For the first machine learning model, they want to start with a simple algorithm, something that gives them
visibility into its working to allow them to validate the usefulness of the problem framing and the data. Simple
models are also easier to implement and deploy, which allow them to quickly build out a framework of exploration. \v

Moving on to the next phase of ``optimizing simple models'', once they have a machine learning framework in place,
they can focus on optimizing the simple machine learning models with different objective functions, hyperparameter
search, feature engineering, more data, and ensembles. \v

Finally, in the last phase of ``complex models'', once they have reached the limit of the simple models and the use
case demands significant model improvement, they start experimenting with more complex models. \v

When selecting a model for your problem, you don't choose from every possible model out there, but usually focus on a
set of models suitable for your problem. When considering what model to use, it's important to consider not only the
model's performance but also its other properties, such as how much data, compute, and time it needs to train, what's
its inference latency, and what's its interpretability. \v

Although there are a lot of thing one needs to be aware of, here are a few tips when searching for models:
\bit
\item Avoid the state-of-the-art trap.
\item Start with the simplest models.
\item Train many quick-and-dirty models from different categories using standard parameters.
\item Measure and compare the performance of the different models.
\item Avoid human biases in selecting models.
\item Evaluate good performance now versus good performance later.
\item Evaluate trade-offs.
\item Understand your model's assumptions.
\eit

\subsubsection{Experiment Tracking \& Versioning}

During the model development process, you often have to experiment with many architectures and many different models
to choose the best one for your problem.

\bd[Experiment]
In machine learning language, an \textbf{experiment} is a systematic procedure to test a hypothesis.
\ed

Some models might seem similar to each other and differ in only one hyperparameter and yet their performances are
dramatically different. In essence any combination that differentiates two runs, is an experiment. In real life, a
typical machine learning engineer runs a lot of experiments before she decides which one is better. It's important to
keep track of all the definitions needed to re-create an experiment and its relevant artifacts.

\bd[Artifact]
An \textbf{artifact} is a file generated during an experiment.
\ed

\be
Examples of artifacts can be files that show the loss curve, evaluation loss graph, logs, or intermediate results of a
model throughout a training process.
\ee

Artifacts enable you to compare different experiments and choose the one best suited for your needs. Comparing
different experiments can also help you understand how small changes affect your model's performance, which, in turn,
gives you more visibility into how your model works. \v

In order to figure out which experiments are worth pursuing, we programmatically track them down as we go and we compare
them.

\bd[Experiment Tracking]
\textbf{Experiment tracking} is the process of tracking the progress and results of an experiment during training.
\ed

Since many problems can arise during the training process, it's important to track what's going on during training not
only to detect and address these issues but also to evaluate whether your model is learning anything useful. In theory,
it's not a bad idea to track everything you can. Most of the time, you probably don't need to look at most of them. But
when something does happen, one or more of them might give you clues to understand and/or debug your model. However,
in practice, due to the limitations of tooling today, it can be overwhelming to track too many things, and tracking
less important things can distract you from tracking what is really important. \v

Following is just a short list of things you might want to consider tracking for each experiment during its training
process:
\bit
\item The loss curve corresponding to the train split and each of the eval splits.
\item The model performance metrics that you care about on all non-test splits.
\item The log of corresponding sample, prediction, and ground truth label.
\item The speed of your model, evaluated by the number of steps per second.
\item System performance metrics such as memory usage and CPU/GPU utilization.
\item The values over time of any parameter and hyperparameter whose changes can affect your model's performance.
\eit

\bd[Versioning]
\textbf{Versioning} is the process of logging all the details of an experiment for the purpose of possibly recreating it
later or comparing it with other experiments.
\ed

Experiment tracking and versioning go hand in hand. There are many tools created for experiment tracking and versioning,
but this gets out of topic for these theoretical notes.

\subsection{Deployment}

Once a final model has been decided, it is a good step to present your solution to gather feedback. Document what you
have done and create a nice presentation, ensuring that your key findings are communicated through beautiful
visualizations or easy -to-remember statements. Make sure you highlight the big picture first, and explain why your
solution achieves the business objective. Don't forget to present interesting points you noticed along the way.
Describe what worked and what did not and most importantly list your assumptions and your system's limitations. \v

Once you present and gather feedback, your model is ready for deployment.

\bd[Deployment]
\textbf{Deployment} is the process of making the model available in production environments in order to provide
predictions to other software systems.
\ed

Deployment is a loose term that generally means making your model running and accessible. When we refer to model
deployment, more often than not we mean deployment of a machine learning pipeline.

\bd[Machine Learning Pipeline]
A \textbf{Machine learning pipeline} is a way to codify and automate the workflow it takes to produce a machine learning
model, consisting of multiple sequential steps that do everything from data extraction and preprocessing to model
training and deployment.
\ed

During model development, your model usually runs in a so-called research, or development, environment.

\bd[Research / Development Environment]
A \textbf{research} or \textbf{development} environment is an isolated environment, without contact to live data, where
research is taken place.
\ed

To be deployed, your model will have to leave the development environment and deployed to a production\footnote{
Production is a spectrum. It might mean generating nice plots in notebooks or keeping your models up and running for
millions of users a day. If you want to deploy a model for a few users to play with, all you have to do is to wrap
your predict function in a POST request endpoint using Flask or FastAPI, put the dependencies this predict function
needs to run in a container, and push your model and its associated container to a cloud service like AWS or GCP to
expose the endpoint. If you want to make your model available to millions of users with a latency of milliseconds and
99\% uptime, set up the infrastructure so that the right person can be immediately notified when something goes wrong,
figure out what went wrong, and seamlessly deploy the updates to fix what's wrong , then you need to do a lot more.}
environment to be used by your end users.

\bd[Production Environment]
A \textbf{production} environment is an environment where the model can receive live data and make predictions that can
be used by other software systems.
\ed

It goes without saying that both research and production machine learning pipelines should be reproducible, i.e.\ when
both research and production environment receive the same raw data, they both should produce the same prediction. \v

In many companies, the responsibility of deploying models falls into the hands of the same people who developed those
models. In many other companies, once a model is ready to be deployed, it will be exported\footnote{Exporting a model
means converting this model into a format that can be used by another application. Some people call this process
``serialization''. There are two parts of a model that you can export: the model definition and the model's parameter
values. The model definition defines the structure of your model, such as how many hidden layers it has and how many
units in each layer. The parameter values provide the values for these units and layers. Usually, these two parts are
exported together.} and handed off to another team to deploy it. However, this separation of responsibilities can
cause high overhead communications across teams and make it slow to update your model. It also can make it hard to
debug should something go wrong. \v

Deploying a model in a production environment includes a cycle between training and inference usually called ``model
development life cycle''.

\bd[Model Development Life Cycle (MDLC)]
\textbf{Model development life cycle} (\textbf{MDLC}) is a term commonly used to describe the flow between training and
inference.
\ed

\fig{mdlc}{0.55}

\subsubsection{Online VS Batch Prediction}

One fundamental decision you'll have to make that will affect both your end users and developers working on your
system is how it generates and serves its predictions to end users: online or batch. The terminologies surrounding
batch and online prediction are still quite confusing due to the lack of standardized practices in the industry.

\bd[Online Prediction]
\textbf{Online prediction} (or \textbf{synchronous prediction}) is when predictions are generated and returned as soon
as requests for these predictions arrive.
\ed

\fig{online}{0.34}

Traditionally, when doing online prediction, requests are sent to the prediction service via REST APIs. \v

A problem with online prediction is that your model might take too long to generate predictions. In that case, there are
three main approaches to reduce its inference latency: make it do inference faster (inference optimization), make the
model smaller (model compression), or make the hardware it's deployed on run faster. \v

Instead of generating predictions as soon as they arrive, what if you compute predictions in advance and store them
in your database, and fetch them when requests arrive? This is exactly what batch prediction does, which is a
workaround for when online prediction isn't cheap enough or isn't fast enough.

\bd[Batch Prediction]
\textbf{Batch prediction} or \textbf{asynchronous prediction} is when predictions are generated periodically or whenever
triggered.
\ed

\fig{batch1}{0.3}

Traditionally, when doing batch prediction, the predictions are stored somewhere, such as in SQL tables or a database,
and retrieved as needed. \v

Batch prediction is good for when you want to generate a lot of predictions and don't need the results immediately.
Because the predictions are precomputed, you don't have to worry about how long it'll take your models to generate
predictions. For this reason, batch prediction can also be seen as a trick to reduce the inference latency of more
complex models—the time it takes to retrieve a prediction is usually less than the time it takes to generate it. \v

The problem with batch prediction is that it makes your model less responsive to users' change preferences. Another
problem with batch prediction is that you need to know what requests to generate predictions for in advance. \v

In many applications, online prediction and batch prediction are used side by side for different use cases. Many
people believe that online prediction is less efficient, both in terms of cost and performance, than batch prediction
because you might not be able to batch inputs together and leverage vectorization or other optimization techniques.
However, this is not necessarily true. Also, with online prediction, you don't have to generate predictions for users
who aren't visiting your site. \v

Here is a table that summarizes the differences between online and batch predictions.

\fig{batchvsonline}{0.45}

As hardware becomes more customized and powerful and better techniques are being developed to allow faster, cheaper
online predictions, online prediction might become the default. In recent years, companies have made significant
investments to move from batch prediction to online prediction. To overcome the latency challenge of online
prediction, two components are required: a real-time pipeline that can work with incoming data and a model that can
generate predictions at a speed acceptable to its end users. For most consumer apps, this means milliseconds.

\subsubsection{On The Cloud VS On The Edge Prediction}

Another decision you'll want to consider is where your model's computation will happen: on the cloud or on the edge.

\bd[On The Cloud]
\textbf{On the cloud} means a large chunk of computation is done on the cloud, either public clouds or private clouds.
\ed

The easiest way is to package your model up and deploy it via a managed cloud service such as AWS or GCP, and this is
how many companies deploy when they get started in machine learning. Cloud services have done an incredible job to
make it easy for companies to bring machine learning models into production. However, there are many downsides to
cloud deployment, with number one being cost, since machine learning models can be compute-intensive, and compute is
expensive.

\bd[On The Edge]
\textbf{On the edge} means a large chunk of computation is done on consumer devices, which are also known as edge
devices.
\ed

Other than help with controlling costs, there are many properties that make edge computing appealing. The first is
that it allows your applications to run where cloud computing cannot. When your models are on public clouds, they
rely on stable internet connections to send data to the cloud and back. Edge computing allows your models to work in
situations where there are no internet connections or where the connections are unreliable. Second, when your models
are already on consumers' devices, you can worry less about network latency. Requiring data transfer over the network
(sending data to the model on the cloud to make predictions then sending predictions back to the users) might make
some use cases impossible. In many cases, network latency is a bigger bottleneck than inference latency. \v

Putting your models on the edge is also appealing when handling sensitive user data. machine learning on the cloud
means that your systems might have to send user data over networks, making it susceptible to being intercepted. Cloud
computing also often means storing data of many users in the same place, which means a breach can affect many people.
Edge computing makes it easier to comply with regulations, like GDPR, about how user data can be transferred or
stored. \v

To move computation to the edge, the edge devices have to be powerful enough to handle the computation, have enough
memory to store machine learning models and load them into memory, as well as have enough battery or be connected to
an energy source to power the application for a reasonable amount of time. Because of the many benefits that edge
computing has over cloud computing, companies are in a race to develop edge devices optimized for different machine
learning use cases.

\subsection{Monitoring}

Software doesn't age like fine wine. It ages poorly. The phenomenon in which a software program degrades over time
even if nothing seems to have changed is known as ``software rot'' or ``bit rot''. Machine learning systems aren't
immune to it. Deploying a model isn't the end of the process. A model's performance degrades over time in production.
Once a model has been deployed, we still have to continually monitor its performance to detect issues as well as
deploy updates to fix them.

\subsubsection{Failures}

Before we identify the cause of machine learning system failures, let's briefly discuss what a machine learning
system failure is. A failure happens when one or more expectations of the system is violated. In traditional
software, we mostly care about a system's operational expectations: whether the system executes its logic within
the expected operational metrics, e.g.\ latency and throughput. For a machine learning system, we care about both its
operational metrics and its machine learning performance metrics. \v

Operational expectation violations are easier to detect, as they're usually accompanied by an operational breakage
such as a timeout, a 404 error on a webpage, an out-of-memory error, or a segmentation fault. However, machine
learning performance expectation violations are harder to detect as doing so requires measuring and monitoring the
performance of machine learning models in production. \v

To effectively detect and fix machine learning system failures in production, it's useful to understand why a model,
after proving to work well during development, would fail in production. We'll examine two types of failures:
``software system failures'' and ``machine learning-specific failures''.

\bd[Software System Failures]
\textbf{Software system failures} are failures that would have happened to non-machine learning systems.
\ed

Keep in mind that just because some failures are not specific to machine learning doesn't mean they're not important
for machine learning engineers to understand. \v

Some common categories of software system failures are:
\bit
\item \textbf{Dependency Failure}: A software package or a codebase that your system depends on breaks, which leads your
system to break. This failure mode is common when the dependency is maintained by a third party, and especially common
if the third party that maintains the dependency no longer exists.
\item \textbf{Deployment Failure}: Failures caused by deployment errors.
\item \textbf{Hardware Failure}: When the hardware that you use to deploy your model doesn't behave the way it should.
\item \textbf{Downtime/Crashing}: If a component of your system runs from a server somewhere and that server is down,
your system will also be down.
\eit

A reason for the prevalence of software system failures is that because machine learning adoption in the industry is
still nascent, tooling around machine learning production is limited and best practices are not yet well developed or
standardized. However, as tooling and best practices for machine learning production mature, there are reasons to
believe that the proportion of software system failures will decrease and the proportion of machine learning-specific
failures will increase.

\bd[Machine Learning-Specific Failures]
\textbf{Machine learning-specific failures} are failures specific to machine learning systems.
\ed

Due to the complexity of machine learning systems and the poor practices in deploying them, a large percentage of
what might look like data shifts on monitoring dashboards are really caused by internal errors. \v

Some common categories of machine learning-specific failures are:
\bit
\item \textbf{Pipeline Differences}: Changes in the training pipeline not correctly replicated in the inference pipeline
and vice versa.
\item \textbf{Training VS Real World Data}: The underlying distribution of the real-world data is unlikely to be the
same as the underlying distribution of the training data. Curating a training dataset that can accurately represent
the data that a model will encounter in production turns out to be very difficult. Real-world data is multifaceted
and, in many cases, virtually infinite, whereas training data is finite and constrained by the time, compute, and
human resources available during the dataset creation and processing.
\item \textbf{Data Distribution Shifts}: The real world isn't stationary. Data distribution shift refers to the
phenomenon when the data a model works with changes over time, which causes this model's predictions to become less
accurate as time passes. Data shifts happen all the time, suddenly, gradually, or seasonally. In general a model does
great when first deployed, but its performance degrades over time as the data distribution changes. This failure mode
needs to be continually monitored and detected for as long as a model remains in production.
\item \textbf{Edge Cases}: Edge cases are the data samples so extreme that they cause the model to make catastrophic
mistakes. Even though edge cases generally refer to data samples drawn from the same distribution, if there is a
sudden increase in the number of data samples in which your model doesn't perform well, it could be an indication
that the underlying data distribution has shifted.
\item \textbf{Degenerate Feedback Loops}: A degenerate feedback loop can happen when the predictions themselves
influence the feedback, which, in turn, influences the next iteration of the model. More formally, a degenerate
feedback loop is created when a system's outputs are used to generate the system's future inputs, which, in turn,
influence the system's future outputs.
\eit

\subsubsection{Monitoring}

As the industry realized that many things can go wrong with a machine learning system, many companies started
investing in monitoring of their machine learning systems.

\bd[Monitoring]
\textbf{Monitoring} refers to the act of tracking, measuring, and logging different metrics that can help us determine
when something goes wrong.
\ed

Monitoring is all about metrics. Because machine learning systems are software systems, the first class of metrics
you'd need to monitor are the operational metrics. These metrics are designed to convey the health of your systems.
They are generally divided into three levels: the network the system is run on, the machine the system is run on, and
the application that the system runs. Examples of these metrics are latency; throughput; the number of prediction
requests your model receives in the last minute, hour, day; the percentage of requests that return with a 2xx code;
CPU/GPU utilization; memory utilization; etc. No matter how good your machine learning model is, if the system is
down, you're not going to benefit from it.

\bd[Service Level Objectives (SLO)]
The conditions to determine whether a system is up are defined in the so called \textbf{service level objectives}
(\textbf{SLOs}).
\ed

\be
For example, an SLO may specify that the service is considered to be up if it has a median latency of less than 200
ms and a 99th percentile under 2 s. A service provider might offer an SLA that specifies their uptime guarantee, such
as 99.99\% of the time, and if this guarantee is not met, they'll give their customers back money.
\ee

Within machine learning-specific metrics, there are generally four artifacts to monitor:
\bit
\item \textbf{Accuracy-Related Metrics}: Accuracy-related metrics are the most direct metrics to help you decide
whether a model's performance has degraded. If your system receives any type of user feedback for the predictions it
makes, it can then be used to calculate your model's accuracy-related metrics.
\item \textbf{Predictions}: Prediction is the most common artifact to monitor. Because each prediction is usually
just a number, predictions are easy to visualize, and their summary statistics are straightforward to compute and
interpret.
\item \textbf{Features}: Monitoring also happens on tracking changes in features, both the features that a model uses
as inputs and the intermediate transformations from raw inputs into final features. Feature monitoring is appealing
because compared to raw input data, features are well structured following a predefined schema. The first step of
feature monitoring is feature validation: ensuring that your features follow an expected schema. The expected schemas
are usually generated from training data or from common sense. If these expectations are violated in production,
there might be a shift in the underlying distribution. \v

There are four major concerns when doing feature monitoring:
\begin{enumerate}
\item A company might have hundreds of models in production, and each model uses hundreds, if not thousands, of
features.
\item While tracking features is useful for debugging purposes, it's not very useful for detecting model performance
degradation.
\item Feature extraction is often done in multiple steps, using multiple libraries, on multiple services.
\item The schema that your features follow can change over time.
\end{enumerate}
\item \textbf{Raw Inputs}: The raw input data might not be easier to monitor, as it can come from multiple sources in
different formats, following multiple structures. The way many machine learning workflows are set up today also makes
it impossible for machine learning engineers to get direct access to raw input data, as the raw input data is often
managed by a data platform team who processes and moves the data to a location like a data warehouse, and the machine
learning engineers can only query for data from that data warehouse where the data is already partially processed.
Therefore, monitoring raw inputs is often a responsibility of the data platform team, not the data science or machine
learning team.
\eit

These four artifacts are generated at four different stages of a machine learning system pipeline, as shown in the
following figure.

\fig{failures}{0.7}

The deeper into the pipeline an artifact is, the more transformations it has gone through, which makes a change in
that artifact more likely to be caused by errors in one of those transformations. However, the more transformations
an artifact has gone through, the more structured it's become and the closer it is to the metrics you actually care
about, which makes it easier to monitor.

\subsubsection{Toolbox}

Measuring, tracking, and interpreting metrics for complex systems is a nontrivial task, and engineers rely on a set
of tools to help them do so. It's common for the industry to herald metrics, logs, and traces as the three pillars of
monitoring. In this section, we will focus on the set of tools from the perspective of monitoring machine learning
systems: logs, dashboards, and alerts.
\bit
\item \textbf{Logs}:Traditional software systems rely on logs to record events produced at runtime. An event is
anything that can be of interest to the system developers, either at the time the event happens or later for
debugging and analysis purposes. When we log an event, we want to make it as easy as possible for us to find it later.
This practice with microservice architecture is called ``distributed tracing''. We want to give each process a
unique ID so that, when something goes wrong, the error message will contain that ID. This allows us to search for
the log messages associated with it. We also want to record with each event all the metadata necessary. Because logs
have grown so large and so difficult to manage, there have been many tools developed to help companies manage and
analyze logs. Analyzing billions of logged events manually is futile, so many companies use machine learning to analyze
logs.
\item \textbf{Dashboards}:Visualizing metrics are critical for monitoring. Since monitoring isn't just for the
developers of a system, but also for non-engineering stakeholders including product managers and business developers,
dashboards is a great tool for making monitoring accessible to non-engineers. Even though graphs can help a lot
with understanding metrics, they aren't sufficient on their own. You still need experience and statistical knowledge.
Graphs are useful for making sense of numbers, but they aren't sufficient. Excessive metrics on a dashboard can
also be counterproductive, a phenomenon known as ``dashboard rot''. It's important to pick the right metrics or
abstract out lower-level metrics to compute higher-level signals that make better sense for your specific tasks.
\item \textbf{Alerts}: When our monitoring system detects something suspicious, it's necessary to alert the right
people about it. An alert consists of the following three components: an alert policy which describes the condition
for an alert, notification channels which describe who is to be notified when the condition is met, and a description
of the alert which helps the alerted person understand what's going on. Depending on the audience of the alert, it's
often necessary to make the alert actionable by providing mitigation instructions or a runbook, a compilation of
routine procedures and operations that might help with handling the alert.
\eit

\subsection{Iteration}

Since a model's performance decays over time, we want to continually update our machine learning model. This process
is widely known as ``iteration''.\footnote{Sometimes iteration is also called ``continual learning'', however this
term might be a bit misleading since continual learning does not mean that a model updates itself with every incoming
sample in production. First this would make your model susceptible to catastrophic forgetting, which is the tendency
of a neural network to completely and abruptly forget previously learned information upon learning new information.
Second, it can make training more expensive, since most hardware backends today were designed for batch processing,
so processing only one sample at a time causes a huge waste of compute power and is unable to exploit data
parallelism. In this notes we will stick to the term iteration.}

\bd[Iteration]
\textbf{Iteration} is the process of continually update a machine learning model in production.
\ed

Iteration combats data distribution shifts, especially when the shifts happen suddenly, and adapts to rare events. \v

The updated model shouldn't be deployed until it's been evaluated. This means that you shouldn't make changes to the
existing model directly. Instead, you create a replica of the existing model and update this replica on new data, and
only replace the existing model with the updated replica if the updated replica proves to be better.

\bd[Champion Model]
An existing model in production is called \textbf{champion model}.
\ed

\bd[Challenger Model]
An updated replica of a champion model is called \textbf{challenger}.
\ed

\fig{replica}{0.5}

There are two types of iterations: ``model iteration'' and ``data iteration''.

\bd[Model Iteration]
In \textbf{model iteration} a new feature is added to an existing model architecture or the model architecture is
changed.
\ed

\bd[Data Iteration]
In \textbf{data iteration} the model architecture and features remain the same, but the model is refreshed with new
data.
\ed

Iteration can happen in two ways: ``stateless retraining'' and ``stateful training''.

\bd[Stateless Retraining]
In \textbf{stateless retraining} the model is trained from scratch each time.
\ed

Stateless retraining is mostly applied for model iteration, as changing the model architecture or adding a new
feature requires training the resulting model from scratch. It usually happens in the early stages of a company,
since in the beginning, the team is focusing on developing new models and updating existing models takes a backseat.
They update an existing model from the scratch when the model's performance has degraded to the point that it's doing
more harm than good, and your team has time to update it.

\bd[Stateful Training]
In \textbf{stateful training}, or \textbf{fine-tuning}, or \textbf{incremental learning}, the model continues training
on new data.
\ed

Stateful training is mostly applied for data iteration, as refreshing the data does not require training the resulting
model from scratch.

\fig{stateless}{0.45}

Since training a model from scratch requires a lot more data than fine-tuning the same model, stateful training
allows you to update your model with less data. Also, with stateful training, it might be possible to avoid storing
data altogether. In the traditional stateless retraining, a data sample might be reused during multiple training
iterations of a model, which means that data needs to be stored. This isn't always possible, especially for data with
strict privacy requirements. In the stateful training paradigm, each model update is trained using only the fresh
data, so a data sample is used only once for training. This means that it's possible to train your model without
having to store data in permanent storage, which helps eliminate many concerns about data privacy. \v

Stateful training doesn't mean no training from scratch since we occasionally train their model from scratch on a
large amount of data to calibrate it. Alternatively, we might also train their model from scratch in parallel with
stateful training and then combine both updated models. Once the infrastructure is set up to allow both stateless
retraining and stateful training, the training frequency is just a knob to twist. You can update your models once an
hour, once a day, or whenever a distribution shift is detected. After all, iteration is about setting up
infrastructure in a way that allows you, to update your models whenever it is needed, whether from scratch or
fine-tuning, and to deploy this update quickly. \v

Even though iteration has many use cases and many companies have applied it with great success, continual
learning still has many challenges:
\bit
\item \textbf{Fresh Data Access Challenge}: The challenge to get fresh data to update your model. Currently, many
companies pull new training data from their data warehouses. The speed at which you can pull data from your data
warehouses depends on the speed at which this data is deposited into your data warehouses. The speed can be slow,
especially if data comes from multiple sources. Being able to pull fresh data isn't enough. If your model needs
labeled data to update, as most models today do, this data will need to be labeled as well. In many applications, the
speed at which a model can be updated is bottlenecked by the speed at which data is labeled. The best candidates for
iteration are tasks where you can get natural labels with short feedback loops.
\item \textbf{Evaluation Challenge}: Making sure that this update is good enough to be deployed. The more frequently
you update your models, the more opportunities there are for updates to fail. Second, iteration makes your
models more susceptible to coordinated manipulation and adversarial attack. Because your models learn online from
real-world data, it makes it easier for users to input malicious data to trick models into learning wrong things.
\eit

These two challenges are usually interconnected. Assuming that the fresher the data, the more likely it is to come
from the current distribution, one idea is to evaluate your model on the most recent data that you have access to. So,
after you've updated your model on the data from the last day, you might want to test this model on the data from the
last hour (assuming that data from the last hour wasn't included in the data used to update your model).

\bd[Backtest]
The method of testing a predictive model on data from a specific period of time in the past is known as a
\textbf{backtest}.
\ed

Because data distributions shift, the fact that a model does well on the data from the last hour doesn't mean that it
will continue doing well on the data in the future. The only way to know whether a model will do well in production
is to deploy it. There are techniques to help you evaluate your models in production safely:
\bit
\item \textbf{Shadow Deployment}: Shadow deployment might be the safest way to deploy a model or any software update.
In shadow deployment we deploy the candidate model in parallel with the existing model. Then for each incoming
request, we route it to both models to make predictions, but only serve the existing model's prediction to the user.
Finally, we log the predictions from the new model for analysis purposes. Only when we've found that the new model's
predictions are satisfactory do we replace the existing model with the new model. Because we don't serve the new
model's predictions to users until you've made sure that the model's predictions are satisfactory, the risk of this
new model doing something funky is low. However, this technique isn't always favorable because it's expensive. It
doubles the number of predictions your system has to generate, which generally means doubling your inference compute
cost.
\item \textbf{A/B Testing\footnote{A/B testing is a way to compare two variants of an object, typically by testing
responses to these two variants, and determining which of the two variants is more effective. In our case, we have
the existing model as one variant, and the candidate model (the recently updated model) as another variant. We'll use
A/B testing to determine which model is better according to some predefined metrics.}}: A/B testing is many machine
learning engineers' first response to how to evaluate machine learning models in production. In an A/B testing we
deploy the candidate model alongside the existing model. Then, a percentage of traffic is routed to the new model for
predictions; the rest is routed to the existing model for predictions. It's common for both variants to serve
prediction traffic at the same time. However, there are cases where one model's predictions might affect another
model's predictions. Finally we monitor and analyze the predictions and user feedback, if any, from both models to
determine whether the difference in the two models' performance is statistically significant.
\item \textbf{Canary Release\footnote{Canary release is a technique to reduce the risk of introducing a new software
version in production by slowly rolling out the change to a small subset of users before rolling it out to the entire
infrastructure and making it available to everybody.}}: In canary release we deploy the candidate model alongside the
existing model. The candidate model is called the canary. Then, a portion of the traffic is routed to the candidate
model. If its performance is satisfactory, we increase the traffic to the candidate model. If not, we abort the
canary and route all the traffic back to the existing model. Finally, we stop when either the canary serves all the
traffic (the candidate model has replaced the existing model) or when the canary is aborted.
\eit

iteration isn't about the retraining frequency, but the manner in which the model is retrained. Once the
infrastructure has been set up to update a model quickly, then the question of how often should we update the models
emerges. To answer that question, we need to figure out how much gain a model will get from being updated with fresh
data. The more gain your model can get from fresher data, the more frequently it should be retrained. In other words,
the question of how often to update a model becomes a lot easier if we know how much the model performance will
improve with updating. One way to figure out the gain is by training your model on the data from different time
windows in the past and evaluating it on the data from today to see how the performance changes.

\subsection{Infrastructure \& MLOps}

So far we discussed the logic for developing machine learning systems, and the considerations for deploying,
monitoring, and continually updating a machine learning system. In this final section we will discuss the underlying
infrastructure and tooling needed to perform all these tasks.

\bd[Infrastructure]
\textbf{Infrastructure} is the set of fundamental facilities and systems that support the sustainable functionality of
households and firms.
\ed

In the machine learning world, infrastructure is the set of fundamental facilities that support the development and
maintenance of machine learning systems. \v

Machine learning systems are complex. The more complex a system, the more it can benefit from good infrastructure.
Infrastructure, when set up right, can help automate processes, reducing the need for specialized knowledge and
engineering time. This, in turn, can speed up the development and delivery of machine learning applications, reduce
the surface area for bugs, and enable new use cases. When set up wrong, however, infrastructure is painful to use and
expensive to replace. \v

It's important to note that every company's infrastructure needs are different, and the infrastructure required for
you depends on the number of applications you develop and how specialized the applications are. Hence, what should be
considered as fundamental facilities, varies greatly from company to company. However there are four different layers
when it comes to infrastructure that every company should pay attention:
\bit
\item \textbf{Storage \& Compute}: The storage layer is where data is collected and stored. The compute layer provides
the compute needed to run your machine learning workloads such as training a model, computing features, generating
features, etc.
\item \textbf{Resource Management}: Resource management comprises tools to schedule and orchestrate your workloads to
make the most out of your available compute resources. Examples of tools in this category include Airflow, Kubeflow,
and Metaflow.
\item \textbf{Machine Learning Platform}: This provides tools to aid the development of machine learning applications
such as model stores, feature stores, and monitoring tools. Examples of tools in this category include SageMaker and
MLflow.
\item \textbf{Development Environment}: This is usually referred to as the dev environment; it is where code is written
and experiments are run. Code needs to be versioned and tested. Experiments need to be tracked.
\eit

Everything we discussed in this section so far, can be included under a common umbrella term called ``Machine Learning 
Operations'' or ``MLOps''.

\bd[Machine Learning Operations (MLOps\footnote{If the definition (or even the name MLOps) sounds familiar, that's
because it pulls heavily from the concept of DevOps, which streamlines the practice of software changes and updates.
Indeed, the two have quite a bit in common like robust automation and trust between teams, the idea of collaboration
and increased communication between teams, the end-to-end service life cycle (build, test, release), and prioritizing
continuous delivery and high quality. Yet there is one critical difference between MLOps and DevOps that makes the
latter not immediately transferable to data science teams: deploying software code into production is fundamentally
different than deploying machine learning models into production. While software code is relatively static, data is
always changing, which means machine learning models are constantly learning and adapting—or not, as the case may be
—to new inputs. The complexity of this environment, including the fact that machine learning models are made up of
both code and data, is what makes MLOps a new and unique discipline.})]
\textbf{Machine learning operations} (\textbf{MLOps}) is a portmanteau of ``machine learning'' and ``operations'' and
it's a set of practices that aims to make machine learning systems production-ready and maintainable.
\ed

MLOps is a term that's been gaining popularity in the industry, quickly becoming a critical component of successful
data science project deployment in the enterprise. It's a process that helps organizations and business leaders
generate long-term value and reduce risk associated with data science, machine learning, and AI initiatives.